\documentclass[11pt]{article}
\usepackage{mathblatt}
% gesetzt von werkzeuge/setzer.py aus dem Lernweg (L3-1 · L3-2 · L3-3 · L3-4 · L3-5 · L3-6 · L3-7) und den Bankzeilen; Muster T6B.
% Präambel des Setzers (werkzeuge/setzer.py), wortgleich aus dem Hand-Satz T6B (bau/proben/2026-10-09/kathete-berechnen), dort aus K4W.
\makeatletter
\setlength{\mbpfbe}{0mm}% keine Punkte (Bauregel 6.4)
% eingerückt (Bauregel 4.7, 6.6): \gEin Einzug, \gSchrift eine Stufe kleiner
\newlength{\gEin}\setlength{\gEin}{0pt}
\let\gSchrift\relax
\renewcommand{\pfkreuzzeile}[1]{\par\noindent\hangindent3.2em\hangafter1\hspace*{1.6em}$\square$\ #1\par}
\newcommand{\gkreuz}[1]{$\square$\ #1\hspace{2em}}
% Bauregel 6.7: links, was man ansieht, rechts, was man tut; Spaltengrenze überall an derselben Stelle
\newsavebox{\gbox}\newlength{\gLinks}
\newcommand{\gzwei}[2]{\par\smallskip\noindent\sbox{\gbox}{#1}%
  \setlength{\gLinks}{\dimexpr0.5\textwidth-\gEin-\mbpfnr\relax}%
  \ifdim\wd\gbox>\dimexpr\gLinks-4mm\relax
    \usebox{\gbox}\par\smallskip\noindent\begin{minipage}{\linewidth}#2\end{minipage}%
  \else
    \begin{minipage}[t]{\gLinks}\vspace{0pt}\usebox{\gbox}\end{minipage}%
    \begin{minipage}[t]{\dimexpr\linewidth-\gLinks\relax}\vspace{0pt}\raggedright #2\end{minipage}%
  \fi\par}
\RenewDocumentEnvironment{pfaufg}{m m m}{%
  \begin{lrbox}{\mbaufgabebox}\begin{minipage}[t]{\linewidth}%
  \gSchrift\noindent\hspace*{\gEin}\mbpfrandmarke{#1}%
  \begin{minipage}[t]{\mbpfnr}\strut\textbf{#2}\end{minipage}%
  \begin{minipage}[t]{\dimexpr\linewidth-\gEin-\mbpfnr\relax}\raggedright\strut\ignorespaces}%
 {\end{minipage}%
  \end{minipage}\end{lrbox}%
  \setlength{\mbaufgabehoehe}{\dimexpr\ht\mbaufgabebox+\dp\mbaufgabebox\relax}%
  \par\addvspace{7pt}\Needspace*{\mbaufgabehoehe}%
  \noindent\usebox{\mbaufgabebox}\par\addvspace{7pt}}
\makeatother
% Rechenraum als Karo (Bauregel 6.9): n Rechenschritte = n mal 10 mm
\newcommand{\karo}[1]{\par\smallskip\noindent\begin{tikzpicture}
  \clip (0,0) rectangle ({\linewidth-1mm},{#1*10mm});
  \draw[black!16,line width=0.3pt,step=5mm] (0,0) grid ({\linewidth},{#1*10mm});
  \end{tikzpicture}\par}
\newcommand{\kk}[1]{$\square$\,#1\hspace{0.9em}}
\newcommand{\linie}[1]{\,{\color{black!45}\rule[-1pt]{#1}{0.4pt}}\,}
% Zeile am Ende jedes Durchgangs: Originale klein grau, darüber; dann Vorgänger/Nachfolger (Bauregel 3.4, 6.5)
\newcommand{\blattende}[2]{\par\vfill\noindent\begin{minipage}{\linewidth}{\scriptsize\color{mbgrau}Originale: #1\par}\smallskip
{\footnotesize #2\par}\end{minipage}}
\newcommand{\pkt}{\textperiodcentered{}}

\newcommand{\kennung}{E5F}
\newcommand{\grau}[1]{{\color{mbgrau}#1}}

\begin{document}
\pfheftstil
\fancyfoot[C]{\parbox[t]{\textwidth}{\mbfhbox\par\vspace{1.5mm}{\footnotesize\color{mbgrau}{\scriptsize \kennung}\hfill\thepage}}}
\pfheftkopf{Prozentwert berechnen}{Klasse 7}
% L3-1: Anschluss an Einheit 1 und 2: Jetzt ist der Prozentsatz gegeben und der graue Teil gesucht; 50 %, 25 %, 10 % sind Hälfte, Viertel, Zehntel des Ganzen
\begin{pfaufg}{}{1.}{}
Der ganze Streifen ist $100\,\%$. Wie viel ist der graue Teil?\par\medskip\noindent
\begin{minipage}[b]{0.32\linewidth}\centering
\begin{tikzpicture}
\fill[black!28] (0,0) rectangle (2.100,0.6);
\draw[black!55,line width=0.4pt] (2.100,0) -- (2.100,0.6);
\draw[line width=0.7pt] (0,0) rectangle (4.2,0.6);
\draw[line width=0.4pt] (0,0.68) -- (0,0.82) -- (2.100,0.82) -- (2.100,0.68);
\node[above,font=\small] at (1.050,0.80) {?};
\node[below,font=\small] at (2.1,0) {ganz: $36$ kg};
\end{tikzpicture}\par\smallskip
\grau{{\small a)\ $50\,\%$ von $36$ kg}\par\smallskip
$36 : 2$\par\smallskip $= 18$ kg}
\end{minipage}\hfill
\begin{minipage}[b]{0.32\linewidth}\centering
\begin{tikzpicture}
\fill[black!28] (0,0) rectangle (1.050,0.6);
\draw[black!55,line width=0.4pt] (1.050,0) -- (1.050,0.6);
\draw[black!55,line width=0.4pt] (2.100,0) -- (2.100,0.6);
\draw[black!55,line width=0.4pt] (3.150,0) -- (3.150,0.6);
\draw[line width=0.7pt] (0,0) rectangle (4.2,0.6);
\draw[line width=0.4pt] (0,0.68) -- (0,0.82) -- (1.050,0.82) -- (1.050,0.68);
\node[above,font=\small] at (0.525,0.80) {?};
\node[below,font=\small] at (2.1,0) {ganz: $80\,€$};
\end{tikzpicture}\par\smallskip
{\small b)\ $25\,\%$ von $80\,€$}\par\smallskip
\linie{30mm}
\end{minipage}\hfill
\begin{minipage}[b]{0.32\linewidth}\centering
\begin{tikzpicture}
\fill[black!28] (0,0) rectangle (0.420,0.6);
\draw[black!55,line width=0.4pt] (0.420,0) -- (0.420,0.6);
\draw[black!55,line width=0.4pt] (0.840,0) -- (0.840,0.6);
\draw[black!55,line width=0.4pt] (1.260,0) -- (1.260,0.6);
\draw[black!55,line width=0.4pt] (1.680,0) -- (1.680,0.6);
\draw[black!55,line width=0.4pt] (2.100,0) -- (2.100,0.6);
\draw[black!55,line width=0.4pt] (2.520,0) -- (2.520,0.6);
\draw[black!55,line width=0.4pt] (2.940,0) -- (2.940,0.6);
\draw[black!55,line width=0.4pt] (3.360,0) -- (3.360,0.6);
\draw[black!55,line width=0.4pt] (3.780,0) -- (3.780,0.6);
\draw[line width=0.7pt] (0,0) rectangle (4.2,0.6);
\draw[line width=0.4pt] (0,0.68) -- (0,0.82) -- (0.420,0.82) -- (0.420,0.68);
\node[above,font=\small] at (0.210,0.80) {?};
\node[below,font=\small] at (2.1,0) {ganz: $450$ m};
\end{tikzpicture}\par\smallskip
{\small c)\ $10\,\%$ von $450$ m}\par\smallskip
\linie{30mm}
\end{minipage}
\fusshilfe{\mbox{1:~b) 20 €; c) 45 m}}
\end{pfaufg}

% L3-2: Aus 10 % lassen sich andere Sätze bauen: mal 3, halbieren, zusammenzählen, von 100 % abziehen
\begin{pfaufg}{}{2.}{}
Rechne zuerst $10\,\%$. Setze daraus zusammen.\par\medskip
\grau{a)\ $30\,\%$ von $80\,€$\quad $10\,\% = 8\,€$, \ $30\,\% = 3 \cdot 8\,€ = 24\,€$}\par\medskip
\noindent\begin{minipage}[t]{0.48\linewidth}
b)\ $20\,\%$ von $45\,€$
\karo{1}
\end{minipage}\hfill
\begin{minipage}[t]{0.48\linewidth}
c)\ $5\,\%$ von $60$ kg
\karo{1}
\end{minipage}
\noindent\begin{minipage}[t]{0.48\linewidth}
d)\ $15\,\%$ von $40\,€$
\karo{1}
\end{minipage}\hfill
\begin{minipage}[t]{0.48\linewidth}
e)\ $35\,\%$ von $200$ m
\karo{1}
\end{minipage}
\noindent\begin{minipage}[t]{0.48\linewidth}
f)\ $90\,\%$ von $30\,€$
\karo{1}
\end{minipage}
\fusshilfe{\mbox{2:~b) 9\,€; c) 3 kg; d) 6\,€; e) 70 m; f) 27\,€}}
\end{pfaufg}

% L3-3: Jeder Satz geht über 1 %: das Ganze durch 100, dann mal den Prozentsatz; warum: Prozent heißt Hundertstel
\begin{pfaufg}{}{3.}{}
Rechne zuerst $1\,\%$, dann mal den Prozentsatz.\par\medskip
\grau{a)\ $7\,\%$ von $300\,€$\quad $1\,\% = 300\,€ : 100 = 3\,€$, \ $7\,\% = 7 \cdot 3\,€ = 21\,€$}\par\medskip
\noindent\begin{minipage}[t]{0.48\linewidth}
b)\ $3\,\%$ von $900$ kg
\karo{1}
\end{minipage}\hfill
\begin{minipage}[t]{0.48\linewidth}
c)\ $12\,\%$ von $400\,€$
\karo{1}
\end{minipage}
\noindent\begin{minipage}[t]{0.48\linewidth}
d)\ $8\,\%$ von $250$ m
\karo{1}
\end{minipage}\hfill
\begin{minipage}[t]{0.48\linewidth}
e)\ $19\,\%$ von $600\,€$
\karo{1}
\end{minipage}
\noindent\begin{minipage}[t]{0.48\linewidth}
f)\ Erkläre, warum du für $1\,\%$ durch $100$ teilst.
\karo{2}
\end{minipage}
\fusshilfe{\mbox{3:~b) 27 kg; c) 48\,€; d) 20 m; e) 114\,€; f) Hundertstel}}
\end{pfaufg}

% L3-4: Kurzweg mit dem Taschenrechner: Prozentsatz als Dezimalzahl mal das Ganze; auch mit Komma im Ganzen und über 100 %
\begin{pfaufg}{}{4.}{}
Schreibe den Prozentsatz als Dezimalzahl und rechne mit dem Taschenrechner.\par\medskip
\grau{a)\ $35\,\%$ von $60\,€$\quad $35\,\% = 0{,}35$; \ $0{,}35 \cdot 60 = 21\,€$}\par\medskip
\noindent\begin{minipage}[t]{0.48\linewidth}
b)\ $64\,\%$ von $25$ kg
\karo{1}
\end{minipage}\hfill
\begin{minipage}[t]{0.48\linewidth}
c)\ $8\,\%$ von $450\,€$
\karo{1}
\end{minipage}
\noindent\begin{minipage}[t]{0.48\linewidth}
d)\ $30\,\%$ von $14{,}60\,€$
\karo{1}
\end{minipage}\hfill
\begin{minipage}[t]{0.48\linewidth}
e)\ $150\,\%$ von $40$ kg
\karo{1}
\end{minipage}
\noindent\begin{minipage}[t]{0.48\linewidth}
f)\ Tom rechnet $40\,\%$ von $75\,€$ so: $0{,}04 \cdot 75 = 3\,€$. Woran siehst du ohne Rechnung, dass das nicht stimmt?
\karo{2}
\end{minipage}
\fusshilfe{\mbox{4:~b) 16 kg; c) 36\,€; d) 4,38\,€; e) 60 kg; f) fast die Hälfte}}
\end{pfaufg}

% L3-5: Erkennen, gemischt mit Einheit 2: Ist ein Teil gesucht (mal) oder ein Prozentsatz (Teil : Ganzes)? Ohne zu rechnen die passende Rechnung wählen
\begin{pfaufg}{}{5.}{}
Kreuze die Rechnung an, die passt. Rechne nicht aus.\par\medskip
\grau{a)\ In einer Klasse sind $25$ Kinder. $40\,\%$ haben ein Haustier. Wie viele Kinder sind das?\par\hspace*{1.2em}$\square$\,$25 : 40$\hspace{0.9em}$\boxtimes$\,$0{,}4 \cdot 25$\hspace{0.9em}$\square$\,$40 : 25$\hspace{0.9em}}\par\medskip
b)\ In einer Klasse sind $25$ Kinder. $10$ fahren mit dem Bus. Wie viel Prozent sind das?\par\hspace*{1.2em}\kk{$0{,}1 \cdot 25$}\kk{$10 : 25$}\kk{$25 : 10$}\par\medskip
c)\ Ein Hemd kostet $40\,€$. Es gibt $15\,\%$ Rabatt. Wie viel Euro spart man?\par\hspace*{1.2em}\kk{$40 : 15$}\kk{$15 : 40$}\kk{$0{,}15 \cdot 40$}\par\medskip
d)\ Eine Hose kostete $50\,€$. Der Rabatt ist $10\,€$. Wie viel Prozent Rabatt sind das?\par\hspace*{1.2em}\kk{$0{,}1 \cdot 50$}\kk{$10 : 50$}\kk{$50 : 10$}\par\medskip
e)\ Eine Schule hat $800$ Schüler. $5\,\%$ fehlen. Wie viele Schüler fehlen?\par\hspace*{1.2em}\kk{$0{,}05 \cdot 800$}\kk{$0{,}5 \cdot 800$}\kk{$800 : 5$}\par\medskip
f)\ Von $800$ Schülern fehlen $40$. Wie viel Prozent sind das?\par\hspace*{1.2em}\kk{$40 : 800$}\kk{$800 : 40$}\kk{$0{,}4 \cdot 800$}\par
\fusshilfe{\mbox{5:~b) $10 : 25$; c) $0{,}15 \cdot 40$; d) $10 : 50$; e) $0{,}05 \cdot 800$; f) $40 : 800$}}
\end{pfaufg}

% L3-6: Die Frage entscheidet: Ersparnis (der Teil) oder neuer Preis (der Rest); am Streifen sind beide zu sehen; in der Prüfung steht der Rest als Falle unter den Antworten
\begin{pfaufg}{}{6.}{}
Ein Handy hat $64$ GB Speicher. $25\,\%$ davon sind mit Fotos belegt.
\gzwei{\begin{tikzpicture}
\fill[black!28] (0,0) rectangle (1.500,0.7);
\draw[black!55,line width=0.4pt] (1.500,0) -- (1.500,0.7);
\draw[black!55,line width=0.4pt] (3.000,0) -- (3.000,0.7);
\draw[black!55,line width=0.4pt] (4.500,0) -- (4.500,0.7);
\draw[line width=0.7pt] (0,0) rectangle (6.0,0.7);
\draw[line width=0.4pt] (0,0.78) -- (0,0.92) -- (1.500,0.92) -- (1.500,0.78);
\node[above,font=\small] at (0.750,0.90) {belegt: $25\,\%$};
\node[below,font=\small] at (3.0,0) {ganz: $64$ GB $= 100\,\%$};
\end{tikzpicture}}{%
\textbf{a)}\ Wie viel GB sind belegt?\linie{14mm}\,GB\par\bigskip
\textbf{b)}\ Wie viel GB sind noch frei?\linie{14mm}\,GB\par\bigskip
\textbf{c)}\ Wie viel Prozent sind frei?\linie{14mm}\,\%}
\fusshilfe{\mbox{6:~a) 16 GB; b) 48 GB; c) 75\,\%}}
\end{pfaufg}

% L3-6: Die Frage entscheidet: Ersparnis (der Teil) oder neuer Preis (der Rest); am Streifen sind beide zu sehen; in der Prüfung steht der Rest als Falle unter den Antworten
\begin{pfaufg}{nach P10 ’21}{7.}{}
Ein Fahrrad kostet $640\,€$. Im Sommer gibt es $15\,\%$ Rabatt. Wie viel Euro spart man? Kreuze an.\par\medskip\noindent\kk{$544\,€$}\kk{$96\,€$}\kk{$9{,}60\,€$}\kk{$42{,}67\,€$}\par
\karo{2}
\fusshilfe{\mbox{7:~96\,€}}
\end{pfaufg}

% L3-7: Zielaufgabe: In einer Prozenttabelle einen Teil ausrechnen, die Zunahme in Prozentpunkten als Anzahl angeben und – Frage gewechselt – einen Prozentsatz bilden, dessen Ganzes erst ausgerechnet wird
\begin{pfaufg}{nach P10 ’19}{8.}{}
Eine Schule hat in beiden Jahren $750$ Schülerinnen und Schüler. Die Tabelle zeigt, wo sie mittags essen.\par\medskip\noindent {\renewcommand{\arraystretch}{1.3}\begin{tabular}{|l|c|c|}\hline
Mittagessen & 2023 & 2025\\\hline Mensa & $36\,\%$ & $48\,\%$\\\hline Pausenbrot & $50\,\%$ & $42\,\%$\\\hline nichts & $14\,\%$ & $10\,\%$\\\hline\end{tabular}}\par\smallskip
\noindent\begin{minipage}[t]{0.48\linewidth}
a)\ Wie viele Schüler aßen 2023 in der Mensa?
\karo{2}
\end{minipage}\hfill
\begin{minipage}[t]{0.48\linewidth}
b)\ Um wie viele Schüler ist die Zahl der Mensa-Esser von 2023 bis 2025 gestiegen?
\karo{3}
\end{minipage}
\noindent\begin{minipage}[t]{0.48\linewidth}
c)\ 2025 essen $72$ der Mensa-Esser vegetarisch. Wie viel Prozent der Mensa-Esser sind das?
\karo{3}
\end{minipage}
\fusshilfe{\mbox{8:~a) 270; b) 90; c) 20\,\%}}
\end{pfaufg}

\par\vfill\noindent{\footnotesize Vorher: Prozentsatz berechnen\quad\textbullet\quad Weiter: Grundwert berechnen\par}
\end{document}
